\begin{frame}{Desigualdad de Van der Corput}

\begin{block}{Lema}
Sean $H,N\in\mathbb N$ con $1\le H\le N$. Entonces

\[
\left|
\sum_{n=1}^{N}
e^{2\pi i f(n)}
\right|^2
\le
\frac{4N}{H}
\sum_{h=0}^{H-1}
\left|
\sum_{n=1}^{N-h}
e^{2\pi i(f(n+h)-f(n))}
\right|.
\]
\end{block}


\begin{alertblock}{Idea}
La suma exponencial asociada a $f(n)$ puede controlarse
mediante sumas exponenciales asociadas a las diferencias

\[
f(n+h)-f(n).
\]
\end{alertblock}

\end{frame}


\begin{frame}{Primer ejemplo: la sucesión $(n^2\alpha)$}

Sea

\[
f(n)=kn^2\alpha.
\]

Entonces

\[
f(n+h)-f(n)
=
k(2hn\alpha+h^2\alpha).
\]


Como $(n\alpha)$ es equidistribuida para
$\alpha\notin\mathbb Q$,

\[
(2hn\alpha+h^2\alpha)
\]

también lo es para cada $h\neq 0$.



Aplicando la desigualdad de Van der Corput y el criterio de Weyl,

\[
(n^2\alpha)
\]

es equidistribuida módulo $1$.

\end{frame}

\begin{frame}{Primer ejemplo: la sucesión $(n^2\alpha)$}

Sea

\[
f(n)=kn^2\alpha.
\]

Entonces

\[
f(n+h)-f(n)
=
k(2hn\alpha+h^2\alpha).
\]

Como $(n\alpha)$ es equidistribuida para
$\alpha\notin\mathbb Q$,

\[
(2hn\alpha+h^2\alpha)
\]

también lo es para cada $h\neq 0$.


Aplicando la desigualdad de Van der Corput y el criterio de Weyl,

\[
(n^2\alpha)
\]

es equidistribuida módulo $1$.

\end{frame}

\begin{frame}{Teorema de las diferencias de Van der Corput}

\begin{block}{Teorema}
Sea $(x_n)$ una sucesión de números reales.

Si para todo entero positivo $h$

\[
(x_{n+h}-x_n)
\]

es equidistribuida módulo $1$,

entonces

\[
(x_n)
\]

también es equidistribuida módulo $1$.
\end{block}


\begin{alertblock}{Importancia}
El resultado permite sustituir una sucesión complicada
por sus sucesiones de diferencias.

Cada aplicación reduce el orden de crecimiento
del problema.
\end{alertblock}

\[
x_n
\longrightarrow
x_{n+h}-x_n
\]

\end{frame}

\begin{frame}{Polinomios con coeficientes irracionales}

\begin{block}{Teorema}
Sea

\[
P(x)=c_mx^m+\cdots+c_1x+c_0.
\]

Si al menos uno de los coeficientes

\[
c_1,\ldots,c_m
\]

es irracional, entonces

\[
(P(n))
\]

es equidistribuida módulo $1$.
\end{block}



\begin{alertblock}{Idea de la demostración}
Las diferencias

\[
P(n+h)-P(n)
\]

son polinomios de grado menor.

Aplicando repetidamente el teorema de las diferencias,
se reduce el problema al caso lineal

\[
n\alpha,
\]

que ya conocemos por el teorema de Weyl.
\end{alertblock}

\end{frame}


\begin{frame}{Teorema de Fejér Generalizado}

\begin{block}{Teorema}
Sea $g:[1,\infty)\to\mathbb R$ una función $k$ veces derivable.

Supongamos que

\[
g^{(k)}(x)\to 0,
\]

\[
x\,|g^{(k)}(x)|\to\infty,
\]

y que $g^{(k)}$ es monótona para $x$
suficientemente grande.

Entonces

\[
(g(n))
\]

es equidistribuida módulo $1$.
\end{block}

\end{frame}

\begin{frame}{Idea de la demostración}

\[
g(n)
\]

\[
\Downarrow
\]

\[
g(n+h)-g(n)
\]

\[
\Downarrow
\]

Derivadas de orden inferior

\[
\Downarrow
\]

Teorema de Fejér
\[
(k=1)
\]


\begin{alertblock}{Idea clave}
El método de Van der Corput permite descender
desde la derivada $k$-ésima hasta el caso tratado
por el teorema de Fejér.
\end{alertblock}

\end{frame}

\begin{frame}{Aplicaciones}

\begin{block}{Consecuencias}
Son equidistribuidas módulo $1$ las sucesiones
\end{block}

\[
(an^\sigma\log^\tau n),
\qquad
0<\sigma<1,
\quad
\tau\in\mathbb R.
\]


\[
(a\log^\tau n),
\qquad
\tau>1.
\]


\[
(an\log^\tau n),
\qquad
\tau<0.
\]


\begin{alertblock}{Conclusión}
El método de Van der Corput extiende el alcance
del teorema de Fejér a familias mucho más amplias
de sucesiones.
\end{alertblock}

\end{frame}